\section*{Chapter \ref{ch:b1:atmospheric-circulation} --- The General Circulation of the Atmosphere}
% values: calc/ch28.py

\begin{solution}{exo:b1:atmospheric-circulation:1}
(a)~$f = 2\Omega\sin\varphi$: $0$ at the equator, \qty{7.29e-5}{s^{-1}} at \ang{30},
\qty{1.031e-4}{s^{-1}} at \ang{45} and \qty{1.458e-4}{s^{-1}} at the pole.
(b)~$fV = \num{1.031e-4} \times 10 = \qty{1.03e-3}{m.s^{-2}}$, about \num{1.05e-4}
of $g$. (c)~In the Northern Hemisphere the deflection is to the right of the motion:
a wind blowing towards the north is pushed towards the east.
\end{solution}

\begin{solution}{exo:b1:atmospheric-circulation:2}
(a)~$|\nabla p|/\rho = (400/\num{4e5})/1.225 = \qty{8.16e-4}{m.s^{-2}}$.
(b)~$10/\num{8.16e-4} = \qty{12250}{s}$, \qty{3.4}{h}. (c)~The Coriolis acceleration
of a wind of \qty{10}{m/s} at \ang{45} N is \qty{1.03e-3}{m.s^{-2}}: the two are of
the same size (ratio $0.79$). Within a few hours of starting to move, the air feels a
Coriolis force able to balance the pressure gradient: the wind settles near the
geostrophic speed, here $\num{8.16e-4}/\num{1.031e-4} = \qty{7.9}{m/s}$, instead of
accelerating down the gradient.
\end{solution}

\begin{solution}{exo:b1:atmospheric-circulation:3}
(a)~$f = 2 \times \num{7.292e-5} \times \sin\ang{40} = \qty{9.374e-5}{s^{-1}}$;
$|\vec v_g| = (400/\num{2.5e5})/(1.225 \times \num{9.374e-5}) = \qty{13.9}{m/s}$, or
$13.93/0.5144 = \qty{27.1}{kn}$. (b)~In the Northern Hemisphere the wind keeps the
low on its left: with the low to the south, it blows towards the west, from the east
(\omterm{def:b1:spherical-astronomy:horizontal}{azimuth} \ang{90}). (c)~The staff points east from the station circle, into the
wind; \qty{27.1}{kn} rounds to \qty{25}{kn}, two full feathers and a half feather,
drawn on the clockwise side of the staff, here its southern side, towards the low.
\end{solution}

\begin{solution}{exo:b1:atmospheric-circulation:4}
(a)~\ang{15} N: north-east \omterm{def:b1:atmospheric-circulation:winds}{trade winds}; \ang{15} S: south-east \omterm{def:b1:atmospheric-circulation:winds}{trade winds};
\ang{45} N and \ang{45} S: \omterm{def:b1:atmospheric-circulation:winds}{westerlies}. (b)~At \ang{15} the surface air of the \omterm{def:b1:atmospheric-circulation:cells}{Hadley
cell} returns towards the equator and is turned to the west by the Coriolis force
(right in the north, left in the south): easterly, negative $u$. At \ang{45} the
surface air of the \omterm{def:b1:atmospheric-circulation:cells}{Ferrel cell} moves poleward and is turned to the east: westerly,
positive $u$. (c)~Of the three \omterm{def:b1:spherical-astronomy:geographic}{latitudes}, the pressure is highest at \ang{30} N,
\qty{1016.8}{hPa}: the \omterm{def:b1:atmospheric-circulation:cells}{subtropical high} lies near \ang{30} N (the full profile puts
its maximum near \ang{35} N).
\end{solution}

\begin{solution}{exo:b1:atmospheric-circulation:5}
(a)~$|\nabla p|/\rho = (200/\num{5e5})/1.225 = \qty{3.27e-4}{m.s^{-2}}$, divided by
$f$: \qty{2.6}{m/s} at \ang{60}, \qty{2.9}{m/s} at \ang{50}, \qty{4.5}{m/s} at
\ang{30}, \qty{6.5}{m/s} at \ang{20}, \qty{12.9}{m/s} at \ang{10} and
\qty{25.7}{m/s} at \ang{5}. (b)~The Coriolis force turns a wind in a time of order
$1/f$: \qty{2.2}{h} at \ang{60}, \qty{11}{h} at \ang{10}, \qty{22}{h} at \ang{5}, and
infinitely long at the equator. Near the equator the pressure field changes, and
friction acts, long before the wind can be turned along the isobars: the balance
never sets in, and the air flows partly down the gradient, at modest speeds.
\end{solution}

\begin{solution}{exo:b1:atmospheric-circulation:6}
(a)~$u = 464.6\,\sin^2\varphi/\cos\varphi$: \qty{14.2}{m/s} at \ang{10},
\qty{32.2}{m/s} at \ang{15}, \qty{57.8}{m/s} at \ang{20}, \qty{134.1}{m/s} at
\ang{30}. (b)~$u = 464.6\,(\cos^2\ang{10} - \cos^2\varphi)/\cos\varphi$:
\qty{17.7}{m/s} at \ang{15}, \qty{42.9}{m/s} at \ang{20}, \qty{117.9}{m/s} at
\ang{30}. (c)~Air that starts farther from the equator starts closer to the axis,
with less angular momentum to spare: it arrives with a weaker westerly at each
\omterm{def:b1:spherical-astronomy:geographic}{latitude}. Spreading the rising branch lowers the predicted jet, though not nearly
enough to reach the observed \qty{40}{m/s}.
\end{solution}

\begin{solution}{exo:b1:atmospheric-circulation:7}
(a)~The wind blows towards the north-east. With one's back to it, facing north-east,
the low is on the left: to the north-west. (b)~The wind blows towards the
south-east; in the Southern Hemisphere the low is on the right of a person with
their back to the wind: to the south-west. (c)~The surface wind is turned from the
isobars towards the low, that is anticlockwise from the \omterm{def:b1:atmospheric-circulation:geostrophic}{geostrophic wind} in the
Northern Hemisphere. The \omterm{def:b1:atmospheric-circulation:geostrophic}{geostrophic wind}, and the low at right angles to its left,
are therefore turned clockwise from the answer of (a): the low lies between
north-west and north, at the \omterm{def:b1:spherical-astronomy:horizontal}{azimuth} $\ang{315} + \alpha$ for a crossing angle
$\alpha$.
\end{solution}

\begin{solution}{exo:b1:atmospheric-circulation:8}
(a)~The area between the equator and \ang{30} N is $2\pi a_\oplus^2\sin\ang{30} =
\qty{1.275e14}{m^2}$; the mass above a square metre between the two levels is
$\Delta p/g = \num{9e4}/9.81 = \qty{9174}{kg}$; the total is
\qty{1.170e18}{kg}. (b)~$\num{1.170e18}/\num{8.7e10} = \qty{1.34e7}{s}$, about
\num{156} days, five months. (c)~$\num{1.170e18}/\num{21.4e10} = \qty{5.47e6}{s}$,
about \num{63} days, two months.
\end{solution}

\begin{solution}{exo:b1:atmospheric-circulation:9}
(a)~$\num{0.01} \times \num{86400} = \qty{864}{m}$ a day. (b)~$0.864 \times 9.8 =
\qty{8.5}{K}$ a day. (c)~The air cools by emitting infrared radiation: the
atmosphere loses more energy by radiation than it absorbs, and the descent of the
air, warming it by compression, makes up the loss. (d)~Sinking air warms and dries
relative to saturation, its clouds evaporate and rain cannot form: the \omterm{def:b1:atmospheric-circulation:cells}{subtropical
highs} are belts of clear skies and little rain, where hot deserts form on land.
\end{solution}

\begin{solution}{exo:b1:atmospheric-circulation:10}
(a)~$V^2 + frV - fr|\vec v_g| = 0$ with $fr = \qty{55.85}{m/s}$:
$V = \tfrac12\bigl(-55.85 + \sqrt{55.85^2 + 4 \times 55.85 \times 15}\bigr) =
\qty{12.3}{m/s}$, \qty{82}{\percent} of the \omterm{def:b1:atmospheric-circulation:geostrophic}{geostrophic wind}. (b)~$V^2 - frV +
fr|\vec v_g| = 0$ has real roots only if $(fr)^2 \geq 4fr|\vec v_g|$, that is
$|\vec v_g| \leq fr/4 = \qty{14.0}{m/s}$. (c)~Near the centre of a high, where $r$ is
small, the limit $fr/4$ is small: the pressure gradient must be weak, the isobars far
apart and the winds light. A depression has no such limit: its isobars can crowd
close to the centre, with strong winds.
\end{solution}

\begin{solution}{exo:b1:atmospheric-circulation:11}
(a)~For small $\varphi$, $\sin\varphi \approx \varphi$ and $\cos\varphi \approx 1$,
so $u \approx \Omega a_\oplus\varphi^2$. (b)~$\ang{10} = \qty{0.1745}{rad}$:
$464.6 \times 0.1745^2 = \qty{14.15}{m/s}$, against \qty{14.23}{m/s} exactly,
\qty{0.5}{\percent} less. (c)~Solving $464.6\sin^2\varphi/\cos\varphi = 27.0$
numerically gives $\varphi = \ang{13.7}$. (d)~The ring already matches the real
annual-mean jet at less than \ang{14}: conservation of angular momentum describes the
upper branch of the \omterm{def:b1:atmospheric-circulation:cells}{Hadley cell} only in the deep tropics; poleward of about
\ang{15}, the eddies take over.
\end{solution}

\begin{solution}{exo:b1:atmospheric-circulation:12}
(a)~Mean $31.7/12 = \ang{2.6}$ N; range from \ang{12.8} S to \ang{10.9} N, that is
\ang{23.7}. (b)~The northernmost position comes in September, three months after the
June solstice; the southernmost in February, two months after the December solstice.
The lag comes from heat storage, above all in the \omterm{def:b1:fluid-envelopes-structure:ocean-layers}{mixed layer} of the ocean
(\cref{ch:b1:fluid-envelopes-structure}), which warms and cools slowly. (c)~The zonal
mean averages all \omterm{def:b1:spherical-astronomy:geographic}{longitudes}: if over the continents the band of rain swings far
into the summer hemisphere, the mean of all \omterm{def:b1:spherical-astronomy:geographic}{longitudes} can cross the equator
although the band over one ocean does not.
\end{solution}

\begin{solution}{pb:b1:atmospheric-circulation:1}
\begin{enumerate}
  \item $\rho = p/(R_dT) = \num{100600}/(287.05 \times 283.15) =
        \qty{1.238}{kg.m^{-3}}$.
  \item $|\nabla p| = \qty{400}{Pa}/\qty{2.5e5}{m} = \qty{1.6e-3}{Pa.m^{-1}}$.
  \item $|\nabla p|/\rho = \num{1.6e-3}/1.238 = \qty{1.293e-3}{m.s^{-2}}$.
  \item $\Omega = 2\pi/\num{86164.2} = \qty{7.292e-5}{rad.s^{-1}}$;
        $f = 2\Omega\sin\ang{50} = \qty{1.117e-4}{s^{-1}}$.
  \item $|\vec v_g| = \num{1.293e-3}/\num{1.117e-4} = \qty{11.57}{m/s}$.
  \item $11.57/0.514444 = \qty{22.5}{kn}$; $11.57 \times 3.6 = \qty{41.7}{km/h}$.
  \item The time is $|\vec v_g|/(|\nabla p|/\rho)$ and $|\vec v_g| = (|\nabla
        p|/\rho)/f$, so it is $1/f = \qty{8951}{s}$, \qty{2.49}{h}; and $2\pi/f =
        \qty{15.6}{h}$. Patterns that last a day or more are therefore close to
        balance.
  \item The \omterm{def:b1:atmospheric-circulation:pgf}{pressure-gradient force} points north, towards the low. The \omterm{def:b1:atmospheric-circulation:geostrophic}{geostrophic
        wind} blows along the isobars with the low on its left: towards the east,
        from the west, \omterm{def:b1:spherical-astronomy:horizontal}{azimuth} \ang{270}: a westerly.
  \item With one's back to a westerly, facing east, the left is north, where the low
        lies: the law is satisfied.
  \item The staff points west from the station circle, towards where the wind comes
        from; \qty{22.5}{kn} rounds to \qty{20}{kn}, two full feathers; they stand on
        the clockwise side of the staff, its northern side, towards the low.
  \item At \ang{50} S, $f < 0$: the wind keeps the low on its right, so it blows
        towards the west, from the east.
  \item $f(\ang{30}) = \qty{7.292e-5}{s^{-1}}$: \qty{17.7}{m/s}; $f(\ang{60}) =
        \qty{1.263e-4}{s^{-1}}$: \qty{10.2}{m/s}.
  \item $f(\ang{5}) = \qty{1.271e-5}{s^{-1}}$ gives \qty{101.7}{m/s}, a wind never
        observed with such gradients. The turning time $1/f$ is \qty{22}{h}:
        the balance has no time to set in, and the formula does not apply.
  \item $9.0/11.57 = 0.778$; along the isobars $9.0\cos\ang{25} = \qty{8.16}{m/s}$,
        across them towards the low $9.0\sin\ang{25} = \qty{3.80}{m/s}$.
  \item The \omterm{def:b1:atmospheric-circulation:geostrophic}{geostrophic wind} blows towards \omterm{def:b1:spherical-astronomy:horizontal}{azimuth} \ang{90}; turned \ang{25} towards
        the low in the north, the surface wind blows towards \ang{65}, from
        \ang{245}, the west-south-west.
  \item $\tan\alpha = 1/(f\tau)$ gives $\tau = 1/(\num{1.117e-4}\tan\ang{25}) =
        \qty{1.92e4}{s}$, \qty{5.3}{h}; the model predicts $11.57\cos\ang{25} =
        \qty{10.5}{m/s}$, more than the observed \qty{9.0}{m/s}: real friction
        is not exactly linear in the wind, and slows it more than the model at this
        angle.
  \item Through a cylinder of radius $R = \qty{500}{km}$ and height $h =
        \qty{1}{km}$: $\dot m = 2\pi R h\rho u_r = 2\pi \times \num{5e5} \times 1000
        \times 1.238 \times 3.80 = \qty{1.48e10}{kg/s}$.
  \item $w = \dot m/(\rho\pi R^2) = 2hu_r/R = 2 \times 1000 \times 3.80/\num{5e5} =
        \qty{0.0152}{m/s}$, about \qty{1310}{m} a day.
  \item $1.31 \times 9.8 = \qty{12.9}{K}$ a day: the air soon reaches saturation,
        and the depression fills with cloud and rain. Under an \omterm{def:b1:atmospheric-circulation:cyclone}{anticyclone} the
        surface air flows outwards, the air above sinks and warms, and the sky
        clears.
  \item $\Omega a_\oplus = \num{7.292e-5} \times \num{6.371e6} = \qty{464.6}{m/s}$
        at the equator; $464.6\cos\ang{30} = \qty{402.3}{m/s}$ at \ang{30} N.
  \item $u = 464.6\,\sin^2\varphi/\cos\varphi$: \qty{57.8}{m/s} at \ang{20} N and
        \qty{134.1}{m/s} at \ang{30} N.
  \item $u = 464.6\,(\cos^2\ang{10} - \cos^2\ang{30})/\cos\ang{30} =
        \qty{117.9}{m/s}$.
  \item $134.1/42.3 = 3.17$. The eddies of the \omterm{def:b1:spherical-astronomy:geographic}{mid-latitudes}, depressions and
        \omterm{def:b1:atmospheric-circulation:cyclone}{anticyclones}, take angular momentum from the upper branch and carry it
        elsewhere; and the air of the cell rises over a band of \omterm{def:b1:spherical-astronomy:geographic}{latitudes} rather
        than at rest on the equator. The zonal mean also averages a meandering jet
        over all \omterm{def:b1:spherical-astronomy:geographic}{longitudes}.
  \item $M/M_0 = (402.3 + 42.3)\cos\ang{30}/464.6 = 0.829$: the air has lost
        \qty{17}{\percent} of the angular momentum it would have kept as a free
        ring.
  \item The \omterm{def:b1:atmospheric-circulation:geostrophic}{geostrophic wind} at P is \qty{11.6}{m/s}, from the west; the wind
        measured at the ground is \qty{9.0}{m/s}, about four fifths of it, and turned
        \ang{25} towards the low by friction.
\end{enumerate}
\end{solution}
